\subsection{Exact extension-field tests for plane quartics}

\begin{claim}{B.Q1}{Proved}{A least-prime-five quartic witness over \(\mathbf F_{3^8}\)}
Let $C/\mathbf F_3$ be the smooth plane quartic
\[
                W^4=X^3Z-XZ^3,\qquad P=(1:0:0).
\]
There is a $P$-split constant cyclic degree-five cover over
$\mathbf F_{3^8}$ whose relative class number is prime to three.
Thus five is the least odd prime different from the characteristic
which gives a witness for this pointed curve over some constant
extension. Every constant degree-five cover of this curve requires
a constant degree divisible by eight; the displayed field is
therefore also the least one for a degree-five witness.
The certificates give six such covers. For each of them the
relative class number stays prime to three exactly under further
constant extensions of odd degree.
\end{claim}
\begin{proof}
The three partial derivatives of the displayed homogeneous
polynomial have no common projective zero: in characteristic
three they are $-Z^3,W^3,X^3$, up to nonzero constants.
The unique point at infinity is $P$. On the affine chart put
$w=W/Z$ and $x=X/Z$. The degree-two morphism
\[
       C\longrightarrow E:\ y^2=x^3-x,
                    \qquad (x,w)\longmapsto(x,w^2)
\]
is branched, in particular at infinity. The elliptic curve has
four points over $\mathbf F_3$, so its Frobenius satisfies
$F^2=[-3]$. Consequently $E[5]$ is constant over
$K=\mathbf F_{3^8}$.

The program \path{code/quartic_cyclic5.py} works in
\[
 K=\mathbf F_3[z]/
       (z^8+2z^5+z^4+2z^2+2z+2).
\]
It specifies a basis of $E[5]$ and all six cyclic subgroups.
For a generator $T=(a,b)$ and $2T=(c,d)$ put
$h=(x-a)(x-c)$. The normalized isogeny has
\[
 r(x)=N(x)/h(x)^2,
 \qquad r'(x)=g(x)/h(x)^3,\qquad
       g=N'h-2Nh',
\]
where
\[
 N=xh^2+
       \sum_{(u,v)\in\{T,2T\}}
        \left(\frac{h^2}{x-u}+\frac{v^2h^2}{(x-u)^2}\right).
\]
These are polynomials. The addition formula on $E$ proves
this isogeny formula by summing $x(R+S)-x(S)$ over the four
nonzero kernel points. The program verifies its defining
identity
\[
 (x^3-x)g^2=N^3+A'Nh^4+B'h^6.
\]
It gives $s,r_0\in K$ with $s^4=-A'$ and
$r_0^3+A'r_0+B'=0$, hence the isogeny to $E$ is
\[
       x_E=(N/h^2-r_0)/s^2,
       y_E=yg/(s^3h^3).
\]
All its kernel points are $K$-rational. Pulling back
$C\to E$ gives the cover whose normalization has equation
\[
       v^4=F(x),\qquad
       F=(x^3-x)g^2h^2/s^6,\qquad
       (X_C,W_C)=((N/h^2-r_0)/s^2,v/h^2).              \tag{B.Q.1}
\]
It is connected: the branched quadratic extension defining
$C/E$ is disjoint from this degree-five unramified isogeny.
Base change makes its map to $C$ finite \emph{\'etale}, cyclic
of degree five, and constant over $K$. Infinity is a rational
point above $P$, and its deck translates split that fibre.
Its genus is $1+5(3-1)=11$.

We give the entire finite cohomological test, not a point-count
guess. Here $\deg F=19$, and $f=x^3-x$, $g$, and $h$ are
pairwise coprime and squarefree, of degrees three, six, and two.
A basis of regular differentials on the normalization is
\[
 \left\{\frac{x^jdx}{v}:0\le j\le3\right\},\quad
 \left\{\frac{gh\,dx}{v^2}\right\},\quad
 \left\{\frac{ghx^jdx}{v^3}:0\le j\le5\right\}.       \tag{B.Q.2}
\]
At a root of $f$ the ramification index is four; at a root of
$gh$ it is two. Checking the orders there proves regularity.
At infinity, $v_\infty(x)=-4$, $v_\infty(v)=-19$ and
$v_\infty(dx)=-5$, giving exactly the displayed ranges.
The eleven forms are independent by their distinct cyclic
characters and powers of $x$, so the genus calculation makes
them a basis.

For a differential $H(x)dx/v^a$, choose $b\in\{1,2,3\}$
with $3b\equiv a\pmod4$. Cartier is computed exactly by
\[
 \frac{1}{v^b}\,\mathcal C_{\mathbf P^1}
       \bigl(H(x)F(x)^{(3b-a)/4}dx\bigr),\qquad
 \mathcal C_{\mathbf P^1}\left(\sum c_ix^i dx\right)
       =\sum c_{3j+2}^{1/3}x^jdx.
\]
Division by $gh$ in the relevant character spaces is checked
with zero remainder. If $M$ is its inverse-Frobenius-semilinear
matrix, eight iterations produce its $K$-linear matrix. Every
entry of $M$, that eighth power, and the cover polynomials is
included in \path{code/quartic_cyclic5.json}. For each of the
six covers the result is
\[
     \det(T-\mathcal C^8)=T^3(T+1)^8.                 \tag{B.Q.3}
\]
The same calculation on $C$, with basis
$dx/w^2,dx/w^3,x\,dx/w^3$, gives Cartier zero.
Serre duality identifies the linear eighth powers of Cartier
and coherent Frobenius as transposes. The Artin--Schreier
sequence identifies the Frobenius-stable coherent subspace
with geometric three-torsion cohomology. Thus the cover's
new three-primary part has rank eight and Frobenius eigenvalue
$-1$ only; its relative class number is a unit. Equation
(B.Q.3) also gives the exact odd-extension condition. This
conclusion is insensitive to possible Jordan blocks.

For completeness, exact enumeration over $\mathbf F_3$,
$\mathbf F_9$, and $\mathbf F_{27}$ gives $4,4,28$ points
on $C$. Newton identities and the reciprocal functional
equation give its Frobenius polynomial
\[
                     (T^2+3)(T^2-3)^2.
\]
For an abelian variety over a finite field,
\citep[Chapter~II, Lemma~1.3 and Remark~1.4]{MilneAV2008}
prove that the algebra generated by Frobenius is a product
of fields. Its minimal polynomial is therefore squarefree,
so $F^4=[9]$ on $J_C$;
equality on a Tate module implies equality of these
endomorphisms. On every order-five character this says
$F^4=-1$ and $F^8=1$. Hence every nonzero character has
period exactly eight, proving the constant-field minimum.
Finally three is excluded in characteristic three and five
is the first admissible odd prime. The independently loaded
certificates verify the finite arithmetic in this argument.
\end{proof}

\begin{claim}{B.Q2}{Proved}{A four-cover all-extension negative test over \(\mathbf F_5\)}
Let $C/\mathbf F_5$ have the same equation and point as in B.Q1.
Pull its quadratic map to $E:y^2=x^3-x$ back along the duals
of the four cyclic degree-three isogenies out of $E$, over
$K=\mathbf F_{5^8}$. These four constant $P$-split covers
have genus seven. Every one has relative class number divisible
by five over $K$ and over every further constant extension.
This assertion concerns these four covers; it does not exclude
the other order-three characters of $J_C$ or other primes.
\end{claim}
\begin{proof}
The program \path{code/quartic_cyclic3_f5.py} computes the four
kernel pairs on $E$ from the roots of its three-division
polynomial $3x^4-6x^2-1$. For a kernel pair generated by
$(a,b)$, the general odd-characteristic isogeny formula is
\[
 h=x-a,\quad N=xh^2+2(3a^2+A)h+4b^2,
 \quad g=N'h-2Nh',
\]
where $E$ has equation $y^2=x^3+Ax+B$. The identity
$fg^2=N^3+A'Nh^4+B'h^6$ is checked coefficient by
coefficient. The image of a different three-torsion point
generates the dual kernel on $E'$. Applying the formula there
gives an isogeny with target coefficients $3^4A,3^6B$;
scaling by three identifies that target with the original $E$.
Thus we have used the dual isogeny, without identifying an
ordinary isogenous elliptic curve with $E$ without proof.

With $f'=x^3+A'x+B'$ on $E'$, the pullback cover is
\[
        v^4=f'g^2h^2/3^6,
        \qquad (X_C,W_C)=(N/(9h^2),v/h^2).
\]
It is a connected constant \emph{\'etale} degree-three
$P$-split cover by the same disjointness and base-change
argument as B.Q1. Now $\deg F=11$, and $\deg(gh)=4$.
The regular basis consists of $x^jdx/v$ for $j=0,1$,
$gh\,dx/v^2$, and $ghx^jdx/v^3$ for $0\le j\le3$.
The same valuation calculation proves it is a full basis.
Cartier extracts the coefficients numbered $5j+4$ from
$H(x)F(x)^a$ in character $a$, and takes fifth roots.

The exact eighth powers have characteristic polynomials
$(T-1)^3$ on the old genus-three cohomology and
$(T-1)^7$ on the genus-seven cover. Polynomial division
therefore gives $(T-1)^4$ for the new factor, with value
zero at one. The old factor has been divided as a polynomial;
we have not divided the two zero values at one.
The new geometric five-torsion has a Frobenius fixed vector.
It remains fixed by every power, proving the all-extension
failure. All four input kernels, both isogenies, the cover
polynomials, and the matrices are in the accompanying JSON.
\end{proof}

\begin{claim}{B.Q3}{Open}{The unbounded quartic classification in B6}
The least witness over arbitrary constant extensions for
\emph{every} smooth plane quartic over $\mathbf F_3$ and
$\mathbf F_5$ remains undecided here. B.Q1 is an exact
least-prime and first-field result for one pointed quartic;
B.Q2 excludes four actual covers at all later fields.
The independent program \path{code/pointed_quartic_tests.py}
also computes the point counts and Weil polynomials of
$X^4+Y^4+Z^4=0$ over $\mathbf F_3$, at $(1:1:1)$,
and $X^4+Y^4-Z^4=0$ over $\mathbf F_5$, at $(1:0:1)$.
Their base-field Jacobian orders are $64$ and $256$,
respectively, so they have no base-field odd-prime Lang
quotient. This is not an all-extension exclusion.

The previous return's large base-field census programs were
not attached and have not been rerun here. We do not treat
their numerical tables as newly verified certificates.
B.CL4 reduces each fixed-prime all-extension question to a
finite calculation, but gives no terminating bound on the
prime search. It also remains undecided whether some curve
has a nonabelian witness of the stated type and has no
class-number witness over any extension. Failure of finitely
many characters, primes, or fields cannot prove that assertion.
\end{claim}
